\subsection{Présentation}
Afin d'accélérer la vitesse d'acquision et apprendre d'avantage sur l'utilisation de Linux en système embarqué,
La proposition suivante utiliserait un système minimal \textit{Buildroot} pour pouvoir compiler et déployer une distribution Linux.

\subsection{PID 1}
Contrairement à \textit{Raspberry PI OS}, le système d'exploitation aurait un programme d'entré minimal sous forme de service.
Celui-ci serait responsable d'attacher les bibliothèques et les pilotes nécessaires pour les sous services approuvé au point 1.
De plus, Ce service doit démarrer les services et s'assurer de redémarrer ceux ayant terminés.
En aucun cas, le \textbf{ServiceManager} doit sortir de la boucle. Autrement, le kernel Linux lancera une erreur et obligera un redémarrage du système.

\subsection{Paquets ajoutés}
\begin{description}
    \item[wpa\_supplicant] Permettera d'utiliser le module WIFI et intéragir avec le pilote
    \item[libgpiod2] Communiquera avec le pilote responsable des \textit{GPIO}
    \item[i2c-dev] Intéragira avec les contrôlleurs I$^2$C pour les moteurs et le capteur de ligne.
    \item[sh] Exécution des scripts \textit{shell} dans le terminal  
    \item[dropbear] Établiera une communication SSH avec l'utilisation du port Ethernet statique ou du WIFI 
\end{description}